\chapter{态矢量、幺正演化与平移对称性}
\providecommand{\ket}[1]{\left\lvert#1\right\rangle}
\providecommand{\bra}[1]{\left\langle#1\right\rvert}

经过两个学季的间隔之后，对量子力学作一次快速回顾是有益的。它不能占据整个学季。前面的主题是量子力学在其最初被发展出来所要解决的问题上的应用：原子、电子及与之相关的粒子，以及光子。

原子先讲。既然电子是原子的组成部分，原子怎么能先于电子呢？起初只需要足够的电子力学来写出氢原子的基本方程。对类电子粒子更详细的研究可以放在后面。在这些应用中，对称性将是一个反复出现的工具：对称性在量子力学中如何被表示，以及它们对一个系统意味着什么。

\section{态、可观测量与正交性}

在经典力学中，一个态由相空间中的一个点表示。在量子力学中，它由复矢量空间中的一个矢量表示。这是一个抽象矢量，而不是普通三维空间中的箭头。态矢量可以相加，也可以乘以复数，而包含它们的空间不必具有任何特定的有限维数。

狄拉克记号把一个态矢量写成一个右矢，
\begin{equation}
\ket{\psi},
\end{equation}
并与之关联一个左矢，
\begin{equation}
\bra{\psi}.
\end{equation}
左矢与右矢通过厄米共轭相联系。把一个左矢和一个右矢放在一起给出一个内积，
\begin{equation}
\langle\phi|\psi\rangle,
\end{equation}
它一般是一个复数。

一个可观测量是可以被测量的东西：位置、动量、能量，或任何其他可测量的量。在量子力学中，一个可观测量由一个线性厄米算符\(A\)表示，
\begin{equation}
A^\dagger=A.
\end{equation}
厄米性是实数性的算符类比。特别地，一个厄米算符的本征值是实数。

\(A\)的一个本征矢量满足
\begin{equation}
A\ket{\alpha}=\alpha\ket{\alpha}.
\end{equation}
可能的本征值就是测量\(A\)的可能结果。例如，如果一个可观测量只能给出\(1\)、\(3\)或\(7\)，那么这些数就是它的本征值。它的本征矢量是使结果确定而非统计性的那些态。在态\(\ket{\alpha}\)中，测量\(A\)确定地给出\(\alpha\)。

内积表达了态之间的一种重要逻辑关系。假设存在一种测量，能够确定地把\(\ket{\phi}\)与\(\ket{\psi}\)区分开来。那么这些态是正交的：
\begin{equation}
\langle\phi|\psi\rangle=0.
\end{equation}
正交性意味着这些态在物理上可以被毫无歧义地区分。对于一个归一化的态，
\begin{equation}
\langle\psi|\psi\rangle=1.
\end{equation}

\section{位置表象中的一个粒子}

首先考虑一个在直线上运动的粒子。它的位置\(x\)是一个可观测量。对于每一个可能的位置\(x_0\)，都存在一个态\(\ket{x_0}\)，在其中位置测量确定地给出\(x_0\)。

与不同位置相关联的态是可区分的，因此是正交的：
\begin{equation}
x\neq x'
\quad\Longrightarrow\quad
\langle x'|x\rangle=0.
\end{equation}
对于一个任意的态\(\ket{\psi}\)，它与一个位置态的重叠定义了波函数：
\begin{equation}
\psi(x)=\langle x|\psi\rangle.
\end{equation}
这就是态矢量的位置表象。

波函数本身并不是一个概率。它可以是正的、负的、虚的，或完全复的。位置概率密度是
\begin{equation}
P(x)=\psi^*(x)\psi(x)=|\psi(x)|^2.
\end{equation}
该乘积是非负的，正如一个概率密度所必须的那样。对于一个归一化的波函数，
\begin{equation}
\int_{-\infty}^{\infty}dx\,|\psi(x)|^2=1.
\end{equation}

在位置表象中，位置算符通过乘法作用：
\begin{equation}
(\hat x\psi)(x)=x\psi(x).
\end{equation}
它的本征函数是集中在一点上的理想化函数。它们被形式地描述为狄拉克{\proseunicodefont δ}函数：以位置本征值为中心的极窄尖峰。它们精确的分布定义对于本回顾并不需要。

在三维中，位置有三个分量，\(x\)、\(y\)和\(z\)，而一个位置态表示三维空间中的一个点。只要两个这样的点在至少一个坐标上不同，它们就是正交的。通常将保留一维记号，而把相应的三维推广理解为不言而喻的。

\section{动量与确定动量态}

对于位置的每一个分量，都有一个相应的动量分量。在位置表象中，沿\(x\)轴的动量算符是
\begin{equation}
\hat p_x=-i\hbar\frac{d}{dx}.
\end{equation}
因此，一个抽象的厄米算符也可以具体地看作对波函数的一种运算。

一个具有确定动量\(p_0\)的态满足本征值方程
\begin{equation}
\hat p_x\psi_{p_0}(x)=p_0\psi_{p_0}(x).
\end{equation}
代入\(\hat p_x\)的微分形式给出
\begin{equation}
-i\hbar\frac{d\psi_{p_0}(x)}{dx}
=p_0\psi_{p_0}(x).
\end{equation}
波函数的导数与波函数本身成正比，所以解具有指数形式：
\begin{equation}
\psi_{p_0}(x)\propto
\exp\left(\frac{ip_0x}{\hbar}\right).
\end{equation}

\begin{classroomqa}
\lecturetimestamp{00:24:26}
\audiencequestion{这个微分方程对任意波函数都成立吗？}
\lecturerresponse{不。它是具有确定动量的态的本征方程。大多数波函数并不满足它；它选出的是动量本征函数。}
\end{classroomqa}

位置本征态与动量本征态之间的区别是鲜明的。位置本征函数集中在一个点上。动量本征函数具有
\begin{equation}
\left|\exp\left(\frac{ip_0x}{\hbar}\right)\right|^2=1.
\end{equation}
因此其模在整个直线上为常数。在这种理想化的意义上，确定的动量意味着位置完全不确定；确定的位置同样意味着动量完全不确定。这是位置--动量不确定关系的极端形式。

\begin{classroomqa}
\lecturetimestamp{00:26:36}
\audiencequestion{如果已知动量为\(p_0\)，为什么仍然存在概率分布？}
\lecturerresponse{所显示的物理量涉及的是位置测量，而不是另一次动量测量。动量与位置是不同的可观测量。对于理想化的动量本征态，位置振幅在整个直线上具有常数的模。动量的振幅属于动量表象，并通过傅里叶变换与位置波函数相联系。}
\end{classroomqa}

\begin{classroomqa}
\lecturetimestamp{00:28:02}
\audiencequestion{当波函数是狄拉克{\proseunicodefont δ}函数时，位置概率是什么样子？}
\lecturerresponse{它集中在指定的位置上。理想化的尖峰表达了位置的确定性。}
\end{classroomqa}

\editorialnote{精确的位置本征态与动量本征态是广义态。狄拉克{\proseunicodefont δ}函数与平面波并不是整条直线上普通的平方可积归一化波函数，而{\proseunicodefont δ}函数的字面平方也不是普通的分布。上述陈述涉及的是它们的理想化局域性与常模。}

波函数可以形式上看作一个向量，对\(x\)的每一个取值都有一个分量。然而，与普通的有限列向量不同，它具有连续无穷多个分量。

\begin{classroomqa}
\lecturetimestamp{00:28:53}
\audiencequestion{态矢量是用列向量表示的吗？}
\lecturerresponse{它们可以那样表示，但位置波函数很难写成普通的列，因为\(x\)取连续多个值。它可以形式上看作一个连续列向量，而相应的左矢则看作一个连续行向量。}
\end{classroomqa}

\section{时间演化与幺正性}

态如何随时间变化？时间演化是态空间上的一个变换。设\(U(t)\)表示经过一段时间\(t\)的演化。那么
\begin{equation}
U(t)\ket{\psi(t_1)}=\ket{\psi(t_1+t)}.
\end{equation}
特别地，
\begin{equation}
\ket{\psi(t)}=U(t)\ket{\psi(0)},
\qquad
U(0)=\mathbf 1.
\end{equation}

对\(U(t)\)的条件来自信息守恒。如果两个初态是可区分的，时间演化绝不能抹去这种区分。正交态必须保持正交。更一般地，任意两个态之间的内积都保持不变：
\begin{equation}
\langle U(t)\phi|U(t)\psi\rangle
=\langle\phi|\psi\rangle.
\end{equation}
由于
\begin{equation}
\langle U(t)\phi|U(t)\psi\rangle
=\langle\phi|U^\dagger(t)U(t)|\psi\rangle,
\end{equation}
所有内积的保持要求
\begin{equation}
U^\dagger(t)U(t)=\mathbf 1.
\end{equation}
满足此方程的算符是幺正的。

\begin{classroomqa}
\lecturetimestamp{00:36:36}
\audiencequestion{保持每一个内积是否比仅仅保持正交性更强？}
\lecturerresponse{对于所讨论的线性态变换，正交性的保持导致内积的保持。普通的旋转提供了熟悉的类比：它们保持正交性与点积。}
\end{classroomqa}

\editorialnote{严格地说，仅保持正交性允许一个整体的非零重标度。物理时间演化已经要求的归一化保持消除了这一自由度，从而给出内积的保持。}

不应将幺正算符与厄米可观测量相混淆。位置与动量是厄米的，但它们一般不是幺正的。幺正算符的本征值是相位，即形如\(e^{i\theta}\)、位于单位圆上的数。

\begin{classroomqa}
\lecturetimestamp{00:38:12}
\audiencequestion{非幺正变换的一个例子是什么？}
\lecturerresponse{位置算符与动量算符就是例子：它们是厄米可观测量，但不是幺正的。幺正算符的本征值具有单位模，因此其本征值是相位。}
\end{classroomqa}

\section{哈密顿量与Schr\"odinger方程}

假定时间演化是连续的。在极短的时间间隔\(\epsilon\)内，演化算符因此必须接近恒等算符：
\begin{equation}
U(\epsilon)=\mathbf 1+\epsilon G+\mathcal O(\epsilon^2).
\end{equation}
利用幺正性，
\begin{align}
U^\dagger(\epsilon)U(\epsilon)
&=
\left(\mathbf 1+\epsilon G^\dagger\right)
\left(\mathbf 1+\epsilon G\right)
+\mathcal O(\epsilon^2)
\\
&=
\mathbf 1+\epsilon(G^\dagger+G)
+\mathcal O(\epsilon^2).
\end{align}
到一阶，
\begin{equation}
G^\dagger=-G.
\end{equation}
因此\(G\)是反厄米的。它可以写成
\begin{equation}
G=-\frac{i}{\hbar}H,
\end{equation}
其中\(H\)是厄米的：
\begin{equation}
H^\dagger=H.
\end{equation}
无穷小时间演化算符因此是
\begin{equation}
U(\epsilon)
=
\mathbf 1-\frac{i\epsilon}{\hbar}H
+\mathcal O(\epsilon^2).
\end{equation}
厄米算符\(H\)就是哈密顿量。

\begin{classroomqa}
\lecturetimestamp{00:45:48}
\audiencequestion{是什么允许\(\epsilon^2\)项被忽略？}
\lecturerresponse{计算是在\(\epsilon\)的一阶下一致地进行的。二阶计算将需要所有二阶项，包括正比于\(-\epsilon^2H^2/(2\hbar^2)\)的项。近似在于\(U(\epsilon)\)的截断表达式，而不在于\(H\)的厄米性。}
\end{classroomqa}

有限时间间隔可以由许多小间隔构建。将\(t\)分成\(N\)个长度为\(t/N\)的步骤。重复无穷小变换给出
\begin{equation}
U(t)
=
\lim_{N\to\infty}
\left(
\mathbf 1-\frac{itH}{N\hbar}
\right)^N
=
\exp\left(-\frac{iHt}{\hbar}\right).
\end{equation}
指数是无穷小变换的有限版本。

现在将态从\(t\)演化到\(t+\epsilon\)：
\begin{align}
\ket{\psi(t+\epsilon)}
&=
U(\epsilon)\ket{\psi(t)}
\\
&=
\left(
\mathbf 1-\frac{i\epsilon}{\hbar}H
\right)\ket{\psi(t)}
+\mathcal O(\epsilon^2).
\end{align}
减去\(\ket{\psi(t)}\)，除以\(\epsilon\)，并取极限，得到
\begin{equation}
\frac{d}{dt}\ket{\psi(t)}
=
-\frac{i}{\hbar}H\ket{\psi(t)}.
\end{equation}
等价地，
\begin{equation}
i\hbar\frac{d}{dt}\ket{\psi(t)}
=
H\ket{\psi(t)}.
\end{equation}
这就是含时Schr\"odinger方程。在位置表象中，它变为
\begin{equation}
i\hbar\frac{\partial}{\partial t}\psi(x,t)
=
H\psi(x,t),
\end{equation}
其中\(H\)作为波函数上的算符作用。

定态Schr\"odinger方程是哈密顿量的本征值方程：
\begin{equation}
H\ket{E}=E\ket{E}.
\end{equation}
本征值\(E\)是可能测量到的能量，即能级。因此，同一个哈密顿量既决定了允许的能量，也决定了态的演化。

如果初始态是能量本征态，它的演化为
\begin{equation}
\ket{E;t}
=
\exp\left(-\frac{iEt}{\hbar}\right)\ket{E}.
\end{equation}
它保持为同一个能量本征矢量，只获得一个相位。\footnote{负号由前面\(U(t)=\exp(-iHt/\hbar)\)的推导固定。后来对相位的一种口头表述具有相反的符号，与该推导不一致。}

\begin{classroomqa}
\lecturetimestamp{01:01:05}
\audiencequestion{能量本征矢量在含时Schr\"odinger方程下如何演化？}
\lecturerresponse{保持右矢不变。由于\(H\ket{E}=E\ket{E}\)，时间演化将能量本征矢量乘以一个由其能量决定的相位：\(\exp(-iEt/\hbar)\)。}
\end{classroomqa}

\section{什么算作对称性？}

时间演化是态空间上的一种幺正变换。对称变换是另一类重要的变换。

旋转对称性对原子物理至关重要。在氢原子中，电子在中心库仑势中运动。旋转坐标系不改变方程，而物理上旋转原子仍然留下一个由相同方程支配的氢原子。态在旋转下可能改变，但动力学定律不改变。

平移给出另一个例子。空空间中的孤立系统在它的每一部分都被位移相同量之后应当表现相同。这个限定很重要。如果存在一个炉子，而只把系统移近它，系统可能表现不同；炉子破坏了该平移对称性。如果系统和炉子一起平移，它们之间的原始关系就恢复了。

并非每个操作都是对称性。将岩盐晶体压缩两倍通常会改变其行为。挤压是一个定义良好的变换，但它一般不是对称性。

因此，对称性是态上的一个保持动力学形式的操作。用\(V\)表示它。它必须保持可区分性：两个正交态在变换后必须仍然正交。相应地，
\begin{equation}
V^\dagger V=\mathbf 1,
\end{equation}
所以\(V\)是幺正的。

\section{对称性与时间演化对易}

假设初始态\(\ket{\psi_1}\)演化为\(\ket{\psi_2}\)：
\begin{equation}
\ket{\psi_2}=U(t)\ket{\psi_1}.
\end{equation}
对一个提议的对称性\(V\)应用于两个态：
\begin{equation}
\ket{\psi_1'}=V\ket{\psi_1},
\qquad
\ket{\psi_2'}=V\ket{\psi_2}.
\end{equation}
如果\(V\)是动力学的对称性，变换后的初始态必须演化为变换后的末态：
\begin{equation}
\ket{\psi_2'}=U(t)\ket{\psi_1'}.
\end{equation}
代入定义给出
\begin{equation}
VU(t)\ket{\psi_1}
=
U(t)V\ket{\psi_1}.
\end{equation}
这必须对每个初始态成立。因此
\begin{equation}
VU(t)=U(t)V.
\end{equation}

两条等价路径可以显示为一个交换图。\footnote{该图重建了在01：10：30和01：16：34之间展开的两路径论证。}
\begin{figure}
\centering
\begin{tikzpicture}[>=Stealth]
\node (a) at (0,0) {$\ket{\psi_1}$};
\node (b) at (4.2,0) {$\ket{\psi_2}=U(t)\ket{\psi_1}$};
\node (c) at (0,-2.0) {$V\ket{\psi_1}$};
\node (d) at (4.2,-2.0) {$V\ket{\psi_2}=U(t)V\ket{\psi_1}$};

\draw[->] (a) -- node[above] {$U(t)$} (b);
\draw[->] (c) -- node[below] {$U(t)$} (d);
\draw[->] (a) -- node[left] {$V$} (c);
\draw[->] (b) -- node[right] {$V$} (d);
\end{tikzpicture}
\caption{对称性操作和时间演化以任一顺序给出相同的末态。}
\end{figure}

\begin{classroomqa}
\lecturetimestamp{01:16:38}
\audiencequestion{对称性必须在所有时刻都与时间演化对易吗？}
\lecturerresponse{是的。动力学的对称性必须对完整的时间演化满足对易条件。}
\end{classroomqa}

\begin{classroomqa}
\lecturetimestamp{01:17:13}
\audiencequestion{这是否意味着可以先旋转再等待，或者先等待再旋转？}
\lecturerresponse{正是如此。如果旋转是一种对称性，那么两种顺序产生相同的变换后末态。}
\end{classroomqa}

对于小时间间隔，
\begin{equation}
U(\epsilon)
=
\mathbf 1-\frac{i\epsilon}{\hbar}H
+\mathcal O(\epsilon^2).
\end{equation}
对易条件变为
\begin{equation}
V\left(
\mathbf 1-\frac{i\epsilon}{\hbar}H
\right)
=
\left(
\mathbf 1-\frac{i\epsilon}{\hbar}H
\right)V
+\mathcal O(\epsilon^2).
\end{equation}
恒等项相消，剩下
\begin{equation}
VH=HV,
\end{equation}
或
\begin{equation}
[V,H]=0.
\end{equation}
因此，一个酉操作是动力学的对称性，当且仅当它与哈密顿量对易。

\begin{classroomqa}
\lecturetimestamp{01:19:29}
\audiencequestion{为什么要求对称性算符\(V\)是酉的？}
\lecturerresponse{对称性必须保持态之间的逻辑区分。它不应把两个互斥的、正交的态变成无法再区分的态。保持这些关系要求酉变换。}
\end{classroomqa}

一个没有显式时间依赖且与\(H\)对易的算符是守恒的。这将对称性与守恒联系起来。在经典力学中，这一关系由诺特定理表述。在当前的量子力学形式中，核心检验是与哈密顿量的对易子。

\section{离散对称性与连续对称性}

对称性可以是离散的或连续的。镜面反射是离散的：它交换左右，两者之间不存在一系列无穷小反射。交换两个全同粒子是另一种离散操作。翻转一个无标记的双面物体提供了同一类例子。

旋转是连续的。一个\(90^\circ\)的旋转可以由许多绕同一轴的小旋转构成。空间平移出于同样的原因是连续的：一个有限位移可以由许多微小位移构成。

这种可分解为小变换的性质赋予连续对称性一种特殊形式。在恒等变换附近，写为
\begin{equation}
V(\epsilon)
=
\mathbf 1-i\epsilon G+\mathcal O(\epsilon^2),
\end{equation}
其中\(\epsilon\)为该变换的参数。酉性意味着
\begin{equation}
G^\dagger=G.
\end{equation}
厄米算符\(G\)称为该对称性的生成元。

如果该变换是一种对称性，那么
\begin{equation}
[H,V(\epsilon)]=0.
\end{equation}
展开到一阶得到
\begin{equation}
[H,G]=0.
\end{equation}
因此，连续对称性由与哈密顿量对易的厄米算符生成。这些生成元是守恒量。

``生成''一词在这里有具体含义。有限变换是通过反复应用无穷小变换构成的，正如有限时间演化是由许多短时间步构成的一样。

\section{平移对称性与动量}

考虑将波函数主动向右平移距离\(\epsilon\)。如果原波函数在\(x_0\)处有峰，则平移后的波函数应在\(x_0+\epsilon\)处有峰。因此，它在坐标\(x\)处的值就是原波函数在\(x-\epsilon\)处的取值：
\begin{equation}
(T(\epsilon)\psi)(x)=\psi(x-\epsilon).
\end{equation}
对于小位移，
\begin{equation}
\psi(x-\epsilon)
=
\psi(x)-\epsilon\frac{d\psi(x)}{dx}
+\mathcal O(\epsilon^2).
\end{equation}
因此
\begin{equation}
T(\epsilon)
=
\mathbf 1-\epsilon\frac{d}{dx}
+\mathcal O(\epsilon^2).
\end{equation}
利用
\begin{equation}
\hat p_x=-i\hbar\frac{d}{dx},
\qquad
\frac{d}{dx}=\frac{i}{\hbar}\hat p_x,
\end{equation}
无穷小平移变为
\begin{equation}
T(\epsilon)
=
\mathbf 1-\frac{i\epsilon}{\hbar}\hat p_x
+\mathcal O(\epsilon^2).
\end{equation}
相应的有限平移为
\begin{equation}
T(\epsilon)
=
\exp\left(-\frac{i\epsilon}{\hbar}\hat p_x\right).
\end{equation}
在约定\(V(\epsilon)=\mathbf 1-i\epsilon G+\cdots\)下，有量纲的平移参数给出
\begin{equation}
G_{\mathrm{trans}}=\frac{\hat p_x}{\hbar}.
\end{equation}
等价地，当\(\hbar\)的因子显示在指数中时，动量被称为空间平移的生成元。

\begin{classroomqa}
\lecturetimestamp{01:34:08}
\audiencequestion{在平移推导中是否丢失了一个负号？}
\lecturerresponse{问题在于赋予平移后波函数的方向。主动向右移动由\(\psi(x-\epsilon)\)表示，而不是\(\psi(x+\epsilon)\)。作出这一修正后，\(T(\epsilon)=\mathbf 1-i\epsilon\hat p_x/\hbar+\cdots\)。}
\end{classroomqa}

动量是否守恒取决于哈密顿量。对于一维自由粒子，
\begin{equation}
H=\frac{\hat p_x^2}{2m}.
\end{equation}
由于每个算符都与其自身的任意函数对易，
\begin{equation}
[H,\hat p_x]=0.
\end{equation}
因此自由粒子哈密顿量在平移下不变，\(\hat p_x\)守恒。从物理上看，无论初始位置选在哪里，粒子的行为都相同。平移初始条件只是平移整个运动。

在三维中，
\begin{equation}
H
=
\frac{\hat p_x^2+\hat p_y^2+\hat p_z^2}{2m}.
\end{equation}
每个动量分量都与\(H\)对易：
\begin{equation}
[H,\hat p_x]=[H,\hat p_y]=[H,\hat p_z]=0.
\end{equation}
每个空间方向都有独立的平移对称性。动量既是守恒量，也是与这些平移相联系的生成元。

下一种情形是旋转对称性。旋转更为复杂，因为绕不同轴的旋转一般不对易：先作\(x\)旋转再作\(y\)旋转，与相反顺序并不相同。这导致了不对易的对称性生成元以及旋转的群论结构。
